% part: intuitionistic-logic
% chapter: introduction
% section: constructive-reasoning

\documentclass[../../../include/open-logic-chapter]{subfiles}

\begin{document}

\olfileid{int}{int}{cr}

\section{રચનાત્મક તર્કવિચાર}

મોડલ સંક્રિયકો કે દ્વિતીય-ક્રમ પરિમાણકો ઉમેરીને પ્રશિષ્ટ તર્કનો
વિસ્તાર કરવાથી ઊલટું, સ્ફુરણાવાદી તર્ક પ્રશિષ્ટ તર્કને મર્યાદિત
કરે છે; એ અર્થમાં તે ``અપ્રશિષ્ટ'' છે. પ્રશિષ્ટ તર્ક અનેક રીતે
\emph{અરચનાત્મક} છે. સ્ફુરણાવાદી તર્કનો આશય રચનાત્મક ગણિતમાં
જોવા મળતા વધુ ``રચનાત્મક'' પ્રકારના તર્કવિચારનું નિદર્શન કરવાનો
છે. નીચેના દાખલા તેની કેટલીક પાયાની પ્રેરણાઓ સ્પષ્ટ કરે છે.

ધારો કે કોઈ દાવો કરે કે તેણે એવી પ્રાકૃતિક સંખ્યા~$n$ નક્કી
કરી છે કે જો $n$ સમ હોય તો રીમાન પરિકલ્પના સત્ય છે, અને જો
$n$ વિષમ હોય તો રીમાન પરિકલ્પના અસત્ય છે. બહુ સારા સમાચાર!{}
રીમાન પરિકલ્પના સત્ય છે કે નહીં એ ગણિતના મોટા વણઉકેલ્યા
પ્રશ્નોમાંનો એક છે; અહીં તો જાણે સમસ્યા ગણતરીના પ્રશ્નમાં,
એટલે કોઈ ચોક્કસ સંખ્યા સમ છે કે નહીં તે નક્કી કરવામાં, ઘટી
ગઈ છે.

તો~$n$નું એ જાદુઈ મૂલ્ય શું છે? તેઓ તેનું વર્ણન આ રીતે કરે છે:
રીમાન પરિકલ્પના સત્ય હોય તો $n$ બરાબર $2$ અને નહીં તો $3$
થાય એવી પ્રાકૃતિક સંખ્યા.

ગુસ્સામાં તમે તમારા પૈસા પાછા માગો છો. પ્રશિષ્ટ દૃષ્ટિએ ઉપરનું
વર્ણન ખરેખર $n$નું એકમાત્ર મૂલ્ય નક્કી કરે છે; પરંતુ તમને
ખરેખર તો \emph{સ્પષ્ટ રીતે} આપેલું $n$નું મૂલ્ય જોઈએ છે.

બીજો, કદાચ થોડો ઓછો કૃત્રિમ, દાખલો જોઈએ. આપણને ખબર છે કે
અસંમેય સંખ્યાને સંમેય ઘાતે ચઢાવવાથી સંમેય પરિણામ મળી શકે છે.
દાખલા તરીકે, $\sqrt{2}^2 = 2$. અસંમેય સંખ્યાને
\emph{અસંમેય} ઘાતે ચઢાવીને સંમેય પરિણામ મેળવી શકાય છે કે
નહીં તે ઓછું સ્પષ્ટ છે. નીચેનો પ્રમેય તેનો હકારાત્મક જવાબ
આપે છે:

\begin{thm}
એવી અસંમેય સંખ્યાઓ $a$ અને $b$ છે કે $a^b$ સંમેય છે.
\end{thm}

\begin{proof}
$\sqrt{2}^{\sqrt{2}}$ને વિચારો. તે સંમેય હોય, તો આપણું કામ
થઈ ગયું: $a = b = \sqrt{2}$ લઈ શકીએ. નહીં તો તે અસંમેય છે.
પછી
\[
(\sqrt{2}^{\sqrt{2}})^{\sqrt{2}} = \sqrt{2}^{\sqrt{2} \cdot
  \sqrt{2}} = \sqrt{2}^2 = 2,
\]
જે સંમેય છે. તેથી આ કિસ્સામાં $a$ને
$\sqrt{2}^{\sqrt{2}}$ અને $b$ને~$\sqrt 2$ લો.
\end{proof}

શું આ પ્રામાણ્ય સાબિતી છે? મોટા ભાગના ગણિતજ્ઞોને એમ લાગે છે.
પણ અહીં ફરી કંઈક થોડું અસંતોષકારક છે: ચોક્કસ ગુણધર્મ ધરાવતા
વાસ્તવિક સંખ્યાઓના જોડાનું અસ્તિત્વ આપણે સાબિત કર્યું છે,
પરંતુ તે સંખ્યાઓનો \emph{કયો} જોડો છે તે કહી શકતા નથી. એ જ
પરિણામ એવી રીતે પણ સાબિત કરી શકાય છે કે જોડો $a$, $b$
સાબિતીમાં જ \emph{આપેલો} હોય: $a = \sqrt{3}$ અને
$b = \log_3 4$ લો. પછી
\[
a^b = \sqrt{3}^{\log_3 4} = 3^{1/2 \cdot \log_3 4} = (3^{\log_3
  4})^{1/2} = 4^{1/2}= 2,
\]
કારણ કે $3^{\log_3 x} = x$.

સ્ફુરણાવાદી તર્કમાં પહેલા પુરાવા જેવી ચાલ માન્ય ન હોય એવા
તર્કવિચારનું નિદર્શન કરવાનો આશય છે. $!A(x)$ સંતોષતા કોઈ~$x$નું
અસ્તિત્વ સાબિત કરવાનો અર્થ, બીજા પુરાવાની જેમ, ચોક્કસ~$x$
અને તે~$!A$ સંતોષે છે તેની સાબિતી આપવાનો છે. $!A$ અથવા
$!B$ સત્ય છે તે સાબિત કરવા માટે બંનેમાંથી એકને સાબિત કરી
શકવું જરૂરી છે.

વિધિવત્ રીતે કહીએ તો, પ્રશિષ્ટ તર્કના !!a{derivation} તંત્રને
ચોક્કસ રીતે મર્યાદિત કરવાથી સ્ફુરણાવાદી તર્ક મળે છે. ગાણિતિક
દૃષ્ટિએ આ માત્ર વિધિવત્ નિષ્પત્તિ-તંત્રો છે; પરંતુ અગાઉ
નોંધ્યું તેમ, તેમનો આશય એક પ્રકારના ગાણિતિક તર્કવિચારનું
નિદર્શન કરવાનો છે. આને ગણિત વિશેના ચોક્કસ તત્ત્વચિંતનાત્મક
દૃષ્ટિકોણથી ન્યાયસંગત તર્કવિચાર (જેમ કે બ્રાઉરનો સ્ફુરણાવાદ),
અથવા વધુ ``મૂર્ત'' અને સંતોષકારક ગાણિતિક તર્કવિચાર (બિશપના
રચનાવાદની દિશામાં) માની શકાય. આ વિધિવત્ વર્ણન અનૌપચારિક
પ્રેરણાને પકડે છે કે નહીં તે અંગે પણ દલીલ થઈ શકે. પરંતુ
આપણું તત્ત્વચિંતનાત્મક વલણ જે હોય તે, સ્ફુરણાવાદી તર્કનો
વિધિવત્ રીતે રજૂ થયેલા તર્ક તરીકે અભ્યાસ કરી શકીએ છીએ;
અને કારણ ગમે તે હોય, ઘણા ગાણિતિક તર્કશાસ્ત્રીઓને તેનો
અભ્યાસ રસપ્રદ લાગે છે.

\end{document}
